\section{Problem 3: Monte Carlo Simulation and Black-Scholes Comparison} \label{c3}
\subsection{Description}
Use the difference equation \eqref{c1:estimation} to generate an ensemble of stock prices
\begin{equation*}
	S^{(k)}_N = S^{(k)}(N\Delta t), k = 1, \dots, M \qquad (\text{where } T = N\Delta t)
\end{equation*}
and then use formula \eqref{c1:monte_carlo_estimation} to compute a Monte Carlo estimate of the value of a five-month call option $\left(T = \dfrac{5}{12} \text{ years} \right)$ for the following parameter values: $r = 0.06$, $\sigma = 0.2$, and $K = \$50$. Find estimates corresponding to current stock prices of $S(0) = s = \$45$, $\$50$, and $\$55$. Use $N = 200$ time steps for each trajectory and $M \approxeq 10{,}000$ sample trajectories for each Monte Carlo estimate. Check the accuracy of your results by comparing the Monte Carlo approximation with the value computed from the exact Black-Scholes formula
\begin{equation}
	C(s) = \dfrac{s}{2} \text{ erfc}\left(-\dfrac{d_1}{\sqrt{2}}\right) - \dfrac{K}{2} e^{-rT} \text{erfc}\left(-\dfrac{d_2}{\sqrt{2}}\right),
	\label{c2:p3:black-scholes}
\end{equation}
where
\[
\begin{cases}
	d_1 = \dfrac{1}{\sigma \sqrt{T}} \left[ \ln\left(\dfrac{s}{k}\right) + \left(r + \dfrac{\sigma^2}{2}\right) T \right] \\[0.3cm]
	d_2 = d_1 - \sigma \sqrt{T}.
\end{cases}
\]
and $\text{erfc}(x)$ is the complementary error function.
\begin{equation*}
	\text{erfc}(x) = \dfrac{2}{\sqrt{\pi}} \int_x^{+\infty} e^{-t^2} \, \mathrm{d}t.
\end{equation*}

\subsection{Simulation and Computed Results}

The Monte Carlo estimation for the value of a European call option is given by \eqref{c1:monte_carlo_estimation},
\begin{equation*}
	\hat{C}(s) = (1 + r \Delta t)^{-N} \left\{ \frac{1}{M} \sum_{k=1}^M f(S_N^{(k)}) \right\},
\end{equation*}

where the stock prices are generated using
\[
	\begin{cases}
		S_{n+1}^{(k)} = S_n^{(k)} + rS_n^{(k)}\Delta t + \sigma S_n^{(k)}\varepsilon_{n+1}^{(k)}\sqrt{\Delta t} \qquad (S_0^{(k)} = s) \\
		f(S) = \max(S - K, 0).
	\end{cases}
\]

For comparison, the exact Black-Scholes value is computed from \eqref{c2:p3:black-scholes},
\begin{equation*}
	C(s) = \dfrac{s}{2} \text{ erfc}\left(-\dfrac{d_1}{\sqrt{2}}\right) - \dfrac{K}{2} e^{-rT} \text{erfc}\left(-\dfrac{d_2}{\sqrt{2}}\right),
\end{equation*}
where
\[
\begin{cases}
	d_1 = \dfrac{1}{\sigma \sqrt{T}} \left[ \ln\left(\dfrac{s}{k}\right) + \left(r + \dfrac{\sigma^2}{2}\right) T \right] \\[0.3cm]
	d_2 = d_1 - \sigma \sqrt{T} \\
	\text{erfc}(x) = \dfrac{2}{\sqrt{\pi}} \displaystyle\int_x^{+\infty} e^{-t^2} \, \mathrm{d}t.
\end{cases}
\]

The simulations were performed with
\begin{equation*}
	\begin{aligned}
		&\text{Parameters:} \quad 
		\begin{cases}
			T = \dfrac{5}{12} \\
			r = 0.06 \\
			\sigma = 0.2 \\
			K = 50
		\end{cases}
		\quad \text{ and } \quad
		&\text{Simulation:} \quad
		\begin{cases}
			N = 200 \\
			M = 10{,}000 \\[0.1cm]
			\Delta t = \dfrac{T}{N} = \dfrac{5/12}{200} = \dfrac{1}{480}.
		\end{cases}
	\end{aligned}
\end{equation*}

The computed results are shown in the table below.

\begin{center}
	\begin{table}
		\centering
		\begin{tabular}{|c|c|c|c|}
			\hline
			\begin{tabular}[c]{@{}c@{}}Initial Stock Price\\ $s$\end{tabular} & 
			\begin{tabular}[c]{@{}c@{}}Monte Carlo Estimate\\ $\hat{C}(s)$\end{tabular} & 
			\begin{tabular}[c]{@{}c@{}}Black-Scholes Value\\ $C(s)$\end{tabular} & 
			Error \\
			\hline
			\$45 & 0.951449 & 0.982198 & 0.030749 \\
			\hline
			\$50 & 3.191711 & 3.206296 & 0.014585 \\
			\hline
			\$55 & 6.817697 & 6.866781 & 0.049084 \\
			\hline
		\end{tabular}
		\caption{Computed results using Monte Carlo and Black-Scholes methods.}
	\end{table}
\end{center}

\subsection{Analysis}
The numerical results show that the Monte Carlo estimates are very close to the corresponding Black-Scholes values for all three initial stock prices. The errors obtained are relatively small, indicating that the simulation method provides accurate approximations of the option prices.

Overall, the results confirm that the Monte Carlo simulation method is capable of producing reliable approximations for European call option pricing.